% !TEX program = xelatex % Указываем, что компилятор - XeLaTeX
\documentclass[aspectratio=169]{beamer}
\usetheme{Madrid}
\usecolortheme{beaver}

% ---------- Encoding & language (xelatex-friendly) ----------
\usepackage{polyglossia} % Вместо babel для XeLaTeX
\setdefaultlanguage{russian}
\setotherlanguages{english}

\usepackage{fontspec} % Для работы со шрифтами
\setmainfont{CMU Serif} % Установка основного шрифта (аналог Computer Modern)
\setsansfont{CMU Sans Serif} % Установка шрифта без засечек
\setmonofont{CMU Typewriter Text} % Установка моноширинного шрифта

% ---------- Common packages ----------
\usepackage{graphicx}
\usepackage{booktabs}
\usepackage{tabularx}
\usepackage{hyperref}
\usepackage{multicol}
\usepackage{ulem}
\usepackage{csquotes}
\usepackage{tikz}
\usetikzlibrary{arrows.meta,positioning,fit,shapes.misc}

% ---------- Listings (Kotlin/Java) ----------
\usepackage{listings} % Используем стандартный listings (inputenc не нужен)
\usepackage[dvipsnames]{xcolor}

\lstdefinelanguage{Kotlin}{
  comment=[l]{//},
  commentstyle={\color{gray}\ttfamily},
  emph={filter, first, firstOrNull, forEach, lazy, map, mapNotNull, println, launch, async, await, withContext, let, run, apply, also, takeIf, takeUnless, repeat, with},
  emphstyle={\color{OrangeRed}},
  identifierstyle=\color{black},
  keywords={!in, !is, abstract, actual, annotation, as, as?, break, by, catch, class, companion, const, constructor, continue, crossinline, data, delegate, do, dynamic, else, enum, expect, external, false, field, file, final, finally, for, fun, get, if, import, in, infix, init, inline, inner, interface, internal, is, lateinit, noinline, null, object, open, operator, out, override, package, param, private, property, protected, public, receiver, reified, return, return@, sealed, set, setparam, super, suspend, tailrec, this, throw, true, try, typealias, typeof, val, var, vararg, when, where, while},
  keywordstyle={\color{NavyBlue}\bfseries},
  morecomment=[s]{/*}{*/},
  morestring=[b]",
  morestring=[s]{"""*}{*"""},
  ndkeywords={@Composable, @Preview, @Deprecated, @JvmField, @JvmName, @JvmOverloads, @JvmStatic, @JvmSynthetic, Array, Byte, Double, Float, Int, Integer, Iterable, Long, Runnable, Short, String, Any, Unit, Nothing, Pair, Triple},
  ndkeywordstyle={\color{BurntOrange}\bfseries},
  sensitive=true,
  stringstyle={\color{ForestGreen}\ttfamily}
}

\lstdefinelanguage{JavaLite}{
  language=Java,
  basicstyle=\footnotesize\ttfamily,
  keywordstyle=\color{NavyBlue}\bfseries,
  commentstyle=\color{gray}\ttfamily,
  stringstyle=\color{ForestGreen}\ttfamily
}

\lstset{
  extendedchars=true,
  basicstyle=\footnotesize\ttfamily,
  breaklines=true,
  showstringspaces=false,
  frame=single,
  framerule=0.2pt,
  tabsize=2,
  numbers=left,
  numberstyle=\tiny,
  numbersep=6pt
}

% ---------- Title ----------
\title{Функциональный стиль в Kotlin}
\subtitle{Лекция 2: Коллекции, цепочки вызовов, безопасная обработка, scope-функции и Pair/Triple}
\author{Александр Глускер}
\institute{Курс «Мобильная разработка»}
\date{\today}

% ---------- TOC at section start ----------
\AtBeginSection[]{
  \begin{frame}{Навигация по лекции}
    \tableofcontents[currentsection, hideallsubsections]
  \end{frame}
}

\begin{document}

% ================== Title ==================
\begin{frame}
  \titlepage
\end{frame}

% ================== Agenda ==================
\begin{frame}{Содержание}
  \begin{columns}[T,onlytextwidth]
    \begin{column}{0.55\textwidth}
      \begin{enumerate}
        \item Что такое функциональный стиль
        \item Основные коллекции: List, Set, Map
        \item Sequence: ленивые вычисления
        \item Методы коллекций: map, filter, reduce и др.
        \item Scope-функции: let, run, apply, takeIf
        \item Безопасная обработка без исключений
        \item Pair и Triple
        \item Практические примеры
      \end{enumerate}
    \end{column}
    \begin{column}{0.45\textwidth}
      \begin{block}{Цели лекции}
        \begin{itemize}
          \item Научиться писать в функциональном стиле
          \item Освоить цепочки вызовов без var и mutable коллекций
          \item Уметь безопасно обрабатывать ошибки
          \item Подготовиться к семинару
        \end{itemize}
      \end{block}
    \end{column}
  \end{columns}
\end{frame}

\section{Что такое функциональный стиль}

\begin{frame}{Функциональный стиль — суть}
  \begin{itemize}
    \item \textbf{Функции как объекты первого класса}: можно передавать, возвращать, хранить
    \item \textbf{Иммутабельность}: данные не изменяются — создаются новые структуры
    \item \textbf{Чистые функции}: одинаковые входы → одинаковые выходы, без побочных эффектов
    \item \textbf{Выражения вместо операторов}: почти всё возвращает значение
    \item \textbf{Декларативность}: описываем \textit{что}, а не \textit{как}
  \end{itemize}
  \begin{block}{В Kotlin}
    Поддержка через лямбды, расширения, стандартную библиотеку коллекций, scope-функции.
  \end{block}
\end{frame}

\begin{frame}{Почему функциональный стиль?}
  \begin{itemize}
    \item Читаемость: цепочки вызовов выражают намерение
    \item Безопасность: меньше ошибок, особенно с null и состоянием
    \item Композируемость: легко комбинировать операции
    \item Тестируемость: чистые функции проще тестировать
    \item Поддержка многопоточности: иммутабельность → нет гонок
  \end{itemize}
  \begin{block}{Важно}
    Не догма — Kotlin позволяет смешивать парадигмы. Используйте функциональный стиль там, где он уместен.
  \end{block}
\end{frame}

\section{Основные коллекции: List, Set, Map}

\begin{frame}[fragile]{List — упорядоченная коллекция}
\lstset{language=Kotlin}
\begin{lstlisting}
val numbers = listOf(1, 2, 3, 4, 5)
val empty = emptyList<Int>()

// Индексация
println(numbers[0]) // 1

// Доступ к первому/последнему
println(numbers.first()) // 1
println(numbers.last())  // 5
\end{lstlisting}
\end{frame}

\begin{frame}[fragile]{Set — уникальные элементы без порядка}
\lstset{language=Kotlin}
\begin{lstlisting}
val unique = setOf(1, 2, 3, 2, 1) // {1, 2, 3}
val emptySet = emptySet<Int>()

// Проверка наличия
println(unique.contains(2)) // true
println(2 in unique)        // true

// Операции над множествами
val a = setOf(1, 2, 3)
val b = setOf(3, 4, 5)
println(a intersect b) // {3}
println(a union b)     // {1, 2, 3, 4, 5}
\end{lstlisting}
\end{frame}

\begin{frame}[fragile]{Map — пары ключ-значение}
\lstset{language=Kotlin}
\begin{lstlisting}
val map = mapOf("a" to 1, "b" to 2, "c" to 3)
val emptyMap = emptyMap<String, Int>()

// Получение значения
println(map["a"])           // 1 (может быть null)
println(map.getValue("a"))  // 1 (бросит исключение, если нет)
println(map.getOrDefault("d", 0)) // 0

// Итерация
map.forEach { (key, value) ->
    println("$key: $value")
}
\end{lstlisting}
\end{frame}

\section{Sequence: ленивые вычисления}

\begin{frame}[fragile]{Зачем Sequence?}
  \textbf{Проблема}: при работе с большими коллекциями, каждый вызов map/filter создает промежуточный список.
  \begin{lstlisting}
  val result = (1..1000000)
      .map { it * 2 }     // Создает список из 1M элементов
      .filter { it > 10 } // Создает еще один список
      .first()
  \end{lstlisting}

  \textbf{Решение}: Sequence — ленивая коллекция. Операции выполняются по одной на элемент.
  \begin{lstlisting}
  val result = (1..1000000).asSequence()
      .map { it * 2 }     // Ничего не вычисляется
      .filter { it > 10 } // Ничего не вычисляется
      .first()            // Вычисляется только до первого подходящего
  \end{lstlisting}
\end{frame}

\begin{frame}[fragile]{Создание и преобразование Sequence}
\lstset{language=Kotlin}
\begin{lstlisting}
// Создание
val seq1 = sequenceOf(1, 2, 3)
val seq2 = (1..10).asSequence()

// Генератор
val fib = generateSequence(Pair(0, 1)) { (a, b) -> Pair(b, a + b) }
    .map { it.first }
    .take(10) // первые 10 чисел Фибоначчи

// Преобразование в коллекцию
val list = seq1.toList()
val set = seq1.toSet()
\end{lstlisting}
\end{frame}

\section{Методы коллекций: map, filter, reduce и др.}

\begin{frame}[fragile]{map и mapNotNull}
\lstset{language=Kotlin}
\begin{lstlisting}
val numbers = listOf(1, 2, 3, 4)

// map: преобразует каждый элемент
val doubled = numbers.map { it * 2 } // [2, 4, 6, 8]

// mapNotNull: отфильтровывает null после преобразования
val strings = listOf("1", "a", "3", "b")
val ints = strings.mapNotNull { it.toIntOrNull() } // [1, 3]
\end{lstlisting}
\end{frame}

\begin{frame}[fragile]{filter, filterNot, filterTo}
\lstset{language=Kotlin}
\begin{lstlisting}
val numbers = listOf(1, 2, 3, 4, 5)

// filter: оставляет, если условие true
val evens = numbers.filter { it % 2 == 0 } // [2, 4]

// filterNot: оставляет, если условие false
val odds = numbers.filterNot { it % 2 == 0 } // [1, 3, 5]

// filterTo: записывает результат в существующую коллекцию
val mutableList = mutableListOf<Int>()
numbers.filterTo(mutableList) { it > 3 } // mutableList = [4, 5]
\end{lstlisting}
\end{frame}

\begin{frame}[fragile]{take, drop, takeWhile, dropWhile}
\lstset{language=Kotlin}
\begin{lstlisting}
val numbers = listOf(1, 2, 3, 4, 5)

// take: первые N элементов
println(numbers.take(3)) // [1, 2, 3]

// drop: пропускает первые N
println(numbers.drop(3)) // [4, 5]

// takeWhile: берет, пока условие true
println(numbers.takeWhile { it < 4 }) // [1, 2, 3]

// dropWhile: пропускает, пока условие true
println(numbers.dropWhile { it < 4 }) // [4, 5]
\end{lstlisting}
\end{frame}

\begin{frame}[fragile]{reduce, fold, sum, maxOrNull}
\lstset{language=Kotlin}
\begin{lstlisting}
val numbers = listOf(1, 2, 3, 4)

// reduce: сворачивает список в одно значение (первый элемент — начальный)
val sum1 = numbers.reduce { acc, n -> acc + n } // 10

// fold: сворачивает с явным начальным значением
val sum2 = numbers.fold(0) { acc, n -> acc + n } // 10

// sum: сумма чисел
val sum3 = numbers.sum() // 10

// maxOrNull: максимальный элемент (или null, если пусто)
val max = numbers.maxOrNull() // 4
\end{lstlisting}
\end{frame}

\begin{frame}[fragile]{groupBy, associate, partition}
\lstset{language=Kotlin}
\begin{lstlisting}
val words = listOf("a", "bb", "ccc", "dd")

// groupBy: группирует по ключу
val byLength = words.groupBy { it.length }
// {1=["a"], 2=["bb", "dd"], 3=["ccc"]}

// associate: создает Map по ключу и значению
val wordToLength = words.associate { it to it.length }
// {a=1, bb=2, ccc=3, dd=2}

// partition: разделяет на две части по условию
val (evens, odds) = numbers.partition { it % 2 == 0 }
// evens = [2, 4], odds = [1, 3, 5]
\end{lstlisting}
\end{frame}

\begin{frame}[fragile]{flatMap, flatten}
\lstset{language=Kotlin}
\begin{lstlisting}
val lists = listOf(listOf(1, 2), listOf(3, 4), listOf(5))

// flatten: объединяет вложенные списки
val flat = lists.flatten() // [1, 2, 3, 4, 5]

// flatMap: преобразует каждый элемент в коллекцию и объединяет
val words = listOf("Hello", "World")
val chars = words.flatMap { it.toList() } // [H, e, l, l, o, W, o, r, l, d]
\end{lstlisting}
\end{frame}

\begin{frame}[fragile]{distinct, sorted, sortedBy}
\lstset{language=Kotlin}
\begin{lstlisting}
val numbers = listOf(3, 1, 4, 1, 5)

// distinct: уникальные элементы
val unique = numbers.distinct() // [3, 1, 4, 5]

// sorted: сортировка по возрастанию
val sorted = numbers.sorted() // [1, 1, 3, 4, 5]

// sortedBy: сортировка по ключу
val words = listOf("bb", "a", "ccc")
val byLength = words.sortedBy { it.length } // [a, bb, ccc]
\end{lstlisting}
\end{frame}

\section{Scope-функции: let, run, apply, takeIf}

\begin{frame}[fragile]{let: безопасная работа с nullable}
\lstset{language=Kotlin}
\begin{lstlisting}
val str: String? = "Hello"

// let выполняется, только если объект не null
str?.let { 
    println(it.length) // 5
    it.uppercase()     // возвращает результат блока
}

// Часто используется для цепочек
val result = str?.let { it.uppercase() } ?: "DEFAULT"
\end{lstlisting}
\end{frame}

\begin{frame}[fragile]{run: выполнить блок и вернуть результат}
\lstset{language=Kotlin}
\begin{lstlisting}
val result = run {
    val x = 10
    val y = 20
    x + y // возвращаемое значение
}
println(result) // 30

// Можно вызывать на объекте
val str = "Hello"
val length = str.run {
    println(this) // "Hello"
    length        // 5
}
\end{lstlisting}
\end{frame}

\begin{frame}[fragile]{apply: конфигурация объекта}
\lstset{language=Kotlin}
\begin{lstlisting}
val list = mutableListOf<String>().apply {
    add("a")
    add("b")
    add("c")
}
// list = ["a", "b", "c"]

// Часто для инициализации
val person = Person().apply {
    name = "Alice"
    age = 30
}
\end{lstlisting}
\end{frame}

\begin{frame}[fragile]{takeIf и takeUnless: условное продолжение}
\lstset{language=Kotlin}
\begin{lstlisting}
val number = 15

// takeIf: возвращает объект, если условие true, иначе null
val evenOrNull = number.takeIf { it % 2 == 0 } // null

// takeUnless: возвращает объект, если условие false
val oddOrNull = number.takeUnless { it % 2 == 0 } // 15

// Комбинируем с let
number.takeIf { it > 10 }
    ?.let { println("Больше 10: $it") } // "Больше 10: 15"
\end{lstlisting}
\end{frame}

\begin{frame}[fragile]{also: побочные эффекты без изменения}
\lstset{language=Kotlin}
\begin{lstlisting}
val list = listOf(1, 2, 3).also {
    println("Список создан: $it") // [1, 2, 3]
}.map { it * 2 }.also {
    println("После map: $it") // [2, 4, 6]
}

// Для логирования в цепочках
val result = computeSomething()
    .also { log("Результат: $it") }
    .process()
\end{lstlisting}
\end{frame}

\section{Безопасная обработка без исключений}

\begin{frame}[fragile]{Цепочка обработки с readlnOrNull()}
\lstset{language=Kotlin}
\begin{lstlisting}
fun main() {
    println("Введите число:")
    val number = readlnOrNull()
        ?.toIntOrNull()              // безопасное преобразование
        ?.takeIf { it > 0 }          // только положительные
        ?.let { "Результат: ${it * 2}" } // преобразование в строку
        ?: "Ошибка: введите положительное число"

    println(number)
}
\end{lstlisting}
\end{frame}

\begin{frame}[fragile]{Обработка списка чисел с ошибками}
\lstset{language=Kotlin}
\begin{lstlisting}
fun main() {
    println("Введите числа через пробел:")
    val sum = readlnOrNull()
        ?.split(" ")                 // разбиваем на части
        ?.mapNotNull { it.toIntOrNull() } // преобразуем, отфильтровываем null
        ?.filter { it > 0 }          // только положительные
        ?.sum()                      // сумма
        ?.takeIf { it < 1000 }       // проверка переполнения
        ?: -1                        // ошибка

    if (sum == -1) {
        println("Ошибка: неверный ввод или сумма >= 1000")
    } else {
        println("Сумма: $sum")
    }
}
\end{lstlisting}
\end{frame}

\section{Pair и Triple}

\begin{frame}[fragile]{Pair: пара значений}
\lstset{language=Kotlin}
\begin{lstlisting}
val pair = Pair("Hello", 42)
val (str, num) = pair // деструктуризация
println(str) // "Hello"
println(num) // 42

// Синтаксический сахар
val pair2 = "Hello" to 42

// Часто используется в map
val map = mapOf("a" to 1, "b" to 2)
\end{lstlisting}
\end{frame}

\begin{frame}[fragile]{Triple: тройка значений}
\lstset{language=Kotlin}
\begin{lstlisting}
val triple = Triple("Alice", 30, "Engineer")
val (name, age, job) = triple
println("$name, $age, $job") // Alice, 30, Engineer

// Пример: возвращаем несколько значений из функции
fun getUserInfo(): Triple<String, Int, String> {
    return Triple("Bob", 25, "Designer")
}
\end{lstlisting}
\end{frame}

\section{Практические примеры}
\section{Практические примеры}

\begin{frame}[fragile]{Пример 1: Безопасный парсер координат}
\lstset{language=Kotlin}
\begin{lstlisting}
fun main() {
    println("Введите: x,y")
    val point = readlnOrNull()
        ?.split(",")
        ?.mapNotNull { it.trim().toDoubleOrNull() }
        ?.takeIf { it.size == 2 }
        ?.let { (x, y) -> "Точка: ($x, $y)" }
        ?: "Ошибка формата"
    println(point)
}
\end{lstlisting}
\end{frame}

\begin{frame}[fragile]{Пример 2: Анализ частотности букв}
\lstset{language=Kotlin}
\begin{lstlisting}
fun main() {
    println("Введите текст:")
    val top3 = readlnOrNull()
        ?.filter { it.isLetter() }
        ?.groupingBy { it.lowercaseChar() }
        ?.eachCount()
        ?.toList()
        ?.sortedByDescending { it.second }
        ?.take(3)
        ?.joinToString { "${it.first}=${it.second}" }
        ?: "Нет данных"
    println("Топ-3 буквы: $top3")
}
\end{lstlisting}
\end{frame}

\begin{frame}[fragile]{Пример 3: Валидация email-ов}
\lstset{language=Kotlin}
\begin{lstlisting}
fun isValidEmail(s: String) = s.count { it == '@' } == 1 &&
    s.split('@').let { it.size == 2 && it[0].isNotEmpty() && '.' in it[1] }

fun main() {
    println("Email'ы через ;")
    val valid = readlnOrNull()
        ?.split(";")
        ?.map { it.trim() }
        ?.filter { isValidEmail(it) }
        ?.ifEmpty { null }
        ?.joinToString("\n")
        ?: "Нет валидных email"
    println(valid)
}
\end{lstlisting}
\end{frame}

\begin{frame}[fragile]{Пример 4: Генератор безопасных паролей}
\lstset{language=Kotlin}
\begin{lstlisting}
fun main() {
    val chars = ('A'..'Z') + ('a'..'z') + ('0'..'9')
    val pwd = (1..12).map { chars.random() }
        .joinToString("")
        .let { p -> if (p.any { it.isDigit() }) p else "Ошибка генерации" }
    println("Пароль: $pwd")
}
\end{lstlisting}
\end{frame}

\begin{frame}[fragile]{Пример 5: Калькулятор выражений}
\lstset{language=Kotlin}
\begin{lstlisting}
fun eval(expr: String): Double? = try {
    val tokens = expr.split(" ")
    tokens.chunked(3).fold(tokens[0].toDouble()) { acc, part ->
        when (part[1]) {
            "+" -> acc + part[2].toDouble()
            "-" -> acc - part[2].toDouble()
            "*" -> acc * part[2].toDouble()
            "/" -> acc / part[2].toDouble()
            else -> return null
        }
    }
} catch (_: Exception) { null }

fun main() {
    println("Введите: a op b op c ...")
    val result = readlnOrNull()?.let(::eval)?.takeIf { it.isFinite() }
    println("Результат: ${result ?: "Ошибка"}")
}
\end{lstlisting}
\end{frame}

\begin{frame}[fragile]{Пример 6: Анализ зависимостей}
\lstset{language=Kotlin}
\begin{lstlisting}
fun main() {
    println("Зависимости: A->B; B->C; ...")
    val deps = readlnOrNull()
        ?.split(";")
        ?.mapNotNull { it.split("->").takeIf { p -> p.size == 2 } }
        ?.associate { it[0].trim() to it[1].trim() }
        ?.let { map ->
            map.keys.filter { it !in map.values }
                .ifEmpty { listOf("Цикл!") }
                .joinToString()
        } ?: "Нет входных данных"
    println("Корневые: $deps")
}
\end{lstlisting}
\end{frame}

\begin{frame}[fragile]{Пример 7: Преобразование единиц измерения}
\lstset{language=Kotlin}
\begin{lstlisting}
val units = mapOf("km" to 1000.0, "m" to 1.0, "cm" to 0.01, "mm" to 0.001)

fun main() {
    println("Введите: значение единица_из -> единица_в")
    val result = readlnOrNull()
        ?.split(" ")
        ?.let { (v, from, _, to) ->
            (v.toDoubleOrNull() ?: return@let null) *
            (units[from] ?: return@let null) /
            (units[to] ?: return@let null)
        }
        ?.takeIf { it.isFinite() && it > 0 }
        ?.let { "%.2f".format(it) }
        ?: "Ошибка конвертации"
    println("Результат: $result")
}
\end{lstlisting}
\end{frame}

\begin{frame}[fragile]{Пример 8: Построение маршрута}
\lstset{language=Kotlin}
\begin{lstlisting}
fun main() {
    println("Маршрут: A-B, B-C, C-D")
    val path = readlnOrNull()
        ?.split(", ")
        ?.map { it.split("-").takeIf { p -> p.size == 2 } }
        ?.filterNotNull()
        ?.fold("") { acc, (from, to) ->
            if (acc.isEmpty() || acc.last() == from[0]) acc + to[0]
            else return@fold "" // разрыв маршрута
        }
        ?.takeIf { it.length > 1 }
        ?: "Невалидный маршрут"
    println("Путь: $path")
}
\end{lstlisting}
\end{frame}
\section{Резюме}

\begin{frame}{Итоги лекции}
  \begin{itemize}
    \item Функциональный стиль — это иммутабельность, чистые функции, выражения
    \item List, Set, Map — основные коллекции; Sequence — для ленивых вычислений
    \item Основные методы: map, filter, reduce, groupBy, sorted, take, drop и др.
    \item Scope-функции: let, run, apply, takeIf — для безопасной и читаемой работы с объектами
    \item Безопасная обработка: цепочки с ?. и ?: вместо try-catch
    \item Pair и Triple — для возврата нескольких значений
    \item Все примеры решены без var и mutable коллекций — как требуется на семинаре
  \end{itemize}
\end{frame}

\begin{frame}{Советы для семинара}
  \begin{itemize}
    \item Начинайте с readlnOrNull() — это ваша точка входа
    \item Используйте ?. для безопасного доступа, ?: для значения по умолчанию
    \item Разбивайте сложные цепочки на логические блоки (даже если пишете в одну строку)
    \item Тестируйте на граничных случаях: пустой ввод, null, ноль, отрицательные числа
    \item Помните: в Kotlin почти всё — выражение. Используйте это!
  \end{itemize}
  \centering
  \Large Удачи на семинаре!
\end{frame}

\end{document}